\section{A finite cross-correspondence criterion for a specified
variety}\label{a-finite-cross-correspondence-criterion-for-a-specified-variety}

\textbf{Scope and input.} Fix \(k=\mathbf F_q\), a prime
\(\ell\ne\operatorname{char}k\), and \(E=\overline{\mathbf Q}_\ell\),
with the algebraic-number embedding used in CC-AG01. This section uses
the proved abelian-sector equivalence CC-AG01, and therefore the
inspected cited all-abelian Tate and numerical--homological inputs
underlying AQ.15. It gives an exact remaining criterion for a specified
smooth projective variety; it does not assert that its cycle condition
always holds. Literature cutoff: 20 November 2026; proof checks: 21
November 2026 UTC.

\subsection{The homological category used
here}\label{the-homological-category-used-here}

Let \(\mathcal H_{\ell,E}(k)\) be the category of pure motives formed
from \textbf{\(E\)-linear algebraic correspondences modulo the kernel of
their \(E\)-valued \(\ell\)-adic cycle class}, followed by idempotent
completion and inversion of the Tate object. Use the usual correction of
commutativity by cohomological parity. Its realization \(R\) on the
ordinary direct sum of cohomology is faithful. Indeed, Poincaré duality
and the Künneth formula identify the class of a correspondence with its
graded linear action; being zero under the action is exactly the
equivalence relation imposed.

This coefficient convention matters. For arbitrary varieties we do not
identify this category with a naive scalar extension of rational
homological Hom spaces without proving scalar injectivity. Nor do we
begin by realizing arbitrary numerical motives. On the abelian--Artin
subcategory, the numerical--homological equality and scalar-injectivity
proof in AQ.15 do identify it with \(\mathcal A_{q,E}\); hence that
subcategory embeds fully faithfully in \(\mathcal H_{\ell,E}(k)\).

All cohomological projectors needed here are algebraic in the
homological category. The characteristic polynomials of Frobenius on
distinct cohomological degrees are relatively prime by their different
weights. The Chinese remainder theorem produces polynomials in the
Frobenius correspondence acting as identity in one degree and zero in
the others. Their idempotent identities hold after realization and thus
in this category. Frobenius semisimplicity is not required for this
construction: the full characteristic polynomials already annihilate the
relevant generalized eigenspaces. This is the finite-field construction
discussed by Milne (1994), author-manuscript pp.~5--6.

\subsection{CC-XC01 --- An abelian target realizing the
semisimplification}\label{cc-xc01-an-abelian-target-realizing-the-semisimplification}

Let \(X/k\) be nonempty, smooth and projective of pure dimension \(d\),
and put

\[
 H=h_{\ell,E}(X),\qquad V=R(H)=\bigoplus_i H^i_{\mathrm{et}}(X_{\bar k},E).
\]

There exist finitely many \(Y_\nu/k\), each an abelian variety or a
finite étale scheme, and an idempotent \(e\) on

\[
 T=\bigoplus_{\nu=1}^s h(Y_\nu),\qquad M=(T,e)\in\mathcal A_{q,E},
\]

such that \(W=R(M)\) is isomorphic, as a Frobenius module, to
\(V^{\mathrm{ss}}\). In particular \(\dim_E W=\dim_E V\). No
semisimplicity assumption on \(V\) is made.

\textbf{Proof.} Deligne's Weil theorem (1974, Theorem 1.6, p.~276) gives
integral algebraic Frobenius eigenvalues of weight \(i\) on \(H^i(X)\).
Their products with their complex conjugates are \(q^i\), so they are
\(p\)-units. Thus every eigenvalue belongs to \(W(q)\), independently of
possible Jordan blocks.

For an eigenvalue \(\alpha\) in a positive degree \(i\), choose
\(\gamma\) with \(\gamma^i=\alpha\). As in AQ-WEIL-EXIST, \(\gamma\) is
an integral weight-one Weil number. Honda's Main Theorem gives an
abelian variety \(B_\alpha/k\) with eigenvalue \(\gamma\) on \(H^1\).
Then \(H^i(B_\alpha^i)\) contains the Künneth summand
\(H^1(B_\alpha)^{\otimes i}\) and hence the eigenvalue \(\alpha\). The
eigenvalues in degree zero are roots of unity, since Frobenius permutes
the finitely many geometric connected components of \(X\). A root of
unity of order \(r\) occurs in
\(H^0(\operatorname{Spec}\mathbf F_{q^r})\).

Use enough copies of these varieties to accommodate every algebraic
multiplicity in the characteristic polynomial of \(V\). This gives a
finite \(T\) whose semisimple realization contains the desired
eigenspaces with at least those multiplicities. Project onto the desired
subspaces. This projector commutes with Frobenius; AQ.15 after
coefficient extension lifts it to \(e\). Its image has exactly
\(V^{\mathrm{ss}}\) as realization. No Tate twist was necessary in this
construction, because all degrees of \(X\) are nonnegative and its
eigenvalues are integral. \(\square\)

\subsection{CC-XC02 --- A determinant condition on one projected
cross-cycle
map}\label{cc-xc02-a-determinant-condition-on-one-projected-cross-cycle-map}

Fix \(T,e,M,W\) as in CC-XC01. A codimension-\(d\) cycle \(z_\nu\) on
\(X\times Y_\nu\) acts in degree zero from \(h(X)\) to \(h(Y_\nu)\).
Define the projected cross-cycle map

\[
 \begin{split}
 \kappa:\ 
 \bigoplus_{\nu=1}^s
 \bigl(\mathrm{CH}^d(X\times Y_\nu)_{\mathbf Q}
              \otimes_{\mathbf Q}E\bigr)
 &\longrightarrow \operatorname{Hom}_{E[F]}(V,W),\\
 (z_\nu)_\nu&\longmapsto R(e)\circ R((z_\nu)_\nu).
 \end{split}                                      \tag{CC-XC.1}
\]

Write \(I=\operatorname{im}\kappa\). It is a finite-dimensional
\(E\)-linear subspace of \(\operatorname{Hom}_E(V,W)\). Choose bases of
the equal-dimensional spaces \(V,W\); the condition
\(\det|_I\not\equiv0\) is independent of those choices.

\textbf{Theorem.} The following conditions are equivalent:

\begin{enumerate}
\def\labelenumi{\arabic{enumi}.}
\tightlist
\item
  \(H\) is isomorphic in \(\mathcal H_{\ell,E}(k)\) to an object of
  \(\mathcal A_{q,E}\).
\item
  \(H\simeq M\) in that category.
\item
  \(I\) contains an invertible linear map.
\item
  The determinant polynomial is not identically zero on \(I\).
\end{enumerate}

When they hold, Frobenius on \(V\) is automatically semisimple. Thus the
determinant condition incorporates both semisimplicity and the needed
cross-correspondence algebraicity.

\textbf{Complete proof.} By the definition of the homological category
and the idempotent completion,

\[
 I=R\bigl(\operatorname{Hom}_{\mathcal H_{\ell,E}}(H,M)\bigr).
                                                        \tag{CC-XC.2}
\]

The codimension in (CC-XC.1) is \(d=\dim X\), independently of
\(\dim Y_\nu\): the action

\[
 a\longmapsto(p_{Y_\nu})_*\bigl(p_X^*a\smile\operatorname{cl}(z_\nu)\bigr)
\]

adds \(2d\) and then subtracts \(2d\) from cohomological degree, with
cancelling Tate twists. The cycle is defined over \(k\), so its action
commutes with Frobenius. Different cohomological degrees have disjoint
Frobenius spectra, so a Frobenius-equivariant map between these
realizations preserves the degrees.

Suppose (3). By (CC-XC.2), choose \(f:H\to M\) with invertible
realization \(\phi:V\to W\). Its Frobenius equivariance shows that \(V\)
is semisimple, since \(W\) is. We must still prove that \(f\) has an
algebraic inverse; faithfulness alone would not establish this.

Poincaré duality gives an isomorphism

\[
 \delta_H:H\xrightarrow{\sim}H^\vee(-d).
\]

Transfer its realization across \(\phi\) to the Frobenius-equivariant
isomorphism

\[
 b=(\phi^\vee(-d))^{-1}\,R(\delta_H)\,\phi^{-1}:
       W\longrightarrow W^\vee(-d).               \tag{CC-XC.3}
\]

Both \(M\) and \(M^\vee(-d)\) belong to \(\mathcal A_{q,E}\). Its full
realization theorem, AQ.15, therefore lifts \(b\) to an algebraic
morphism \(\delta_M:M\to M^\vee(-d)\). This step requires fullness only
inside the already established abelian sector.

Now form the algebraic morphism

\[
 g=\delta_H^{-1}\circ f^\vee(-d)\circ\delta_M:
          M\longrightarrow H.                    \tag{CC-XC.4}
\]

Using (CC-XC.3), its realization is

\[
 R(g)=R(\delta_H)^{-1}\,
       \phi^\vee(-d)\,
       (\phi^\vee(-d))^{-1}\,
       R(\delta_H)\,\phi^{-1}
       =\phi^{-1}.
\]

Faithfulness of the homological realization now gives \(gf=1_H\) and
\(fg=1_M\). Thus (3) implies (2). This argument needs neither a separate
surjectivity assumption for reverse cross-correspondences nor
numerical--homological equivalence for \(X\).

Condition (2) implies (1), and an isomorphism in (2) has invertible
realization, so it implies (3). If (1) holds, say \(H\simeq M'\) with
\(M'\in\mathcal A_{q,E}\), then \(V\) is semisimple. Hence
\(R(M')\simeq V\simeq V^{\mathrm{ss}}\simeq W\). Full faithfulness
within \(\mathcal A_{q,E}\) lifts this Frobenius-module isomorphism and
its inverse to \(M'\simeq M\); therefore (1) implies (2).

Finally (3) implies (4). Conversely, a nonzero polynomial on a
finite-dimensional vector space over the infinite field \(E\) is nonzero
at some point: after choosing coordinates this follows by induction on
the number of variables, using the finitely many roots of a nonzero
one-variable polynomial. Apply this to \(\det|_I\) to deduce (3) from
(4). This proves all equivalences and the semisimplicity conclusion.
\(\square\)

\textbf{Finite witnesses.} The image \(I\) is spanned by the images of
algebraic cycles and is finite-dimensional. Consequently it has a finite
spanning collection of such images. Equivalently, the condition is
witnessed by finitely many cycles \(z_1,\ldots,z_t\) for which

\[
 \det\!\left(\sum_{j=1}^t x_j\kappa(z_j)\right)
\]

is not the zero polynomial. This is an existence criterion; it supplies
no effective method or complexity bound for finding the cycles.

\subsection{CC-XC03 --- The semisimple version as a finite projected
Tate
condition}\label{cc-xc03-the-semisimple-version-as-a-finite-projected-tate-condition}

\textbf{Corollary.} If Frobenius on \(V\) is semisimple, then the
conditions in CC-XC02 are also equivalent to the surjectivity of
\(\kappa\).

\textbf{Proof.} Under semisimplicity, \(V\simeq W\) as Frobenius
modules. Surjectivity then puts such an isomorphism in \(I\), so CC-XC02
applies. Conversely, if \(H\simeq M\), transport any
Frobenius-equivariant map \(V\to W\) across that isomorphism to an
endomorphism of \(W\). Full faithfulness on \(M\) supplies its algebraic
endomorphism. Composing gives the required morphism \(H\to M\), and
(CC-XC.2) places its realization in \(I\). Thus \(\kappa\) is onto.
\(\square\)

One sufficient hypothesis for this projected surjectivity is the usual
Tate surjectivity in \textbf{codimension \(d\)} on each of the finitely
many \(X\times Y_\nu\):

\[
 \mathrm{CH}^d(X\times Y_\nu)_{\mathbf Q}
       \otimes_{\mathbf Q}\mathbf Q_\ell
 \twoheadrightarrow
 H^{2d}((X\times Y_\nu)_{\bar k},\mathbf Q_\ell(d))^{F=1}.
                                                        \tag{CC-XC.5}
\]

Indeed, Poincaré duality and Künneth identify this invariant space with
the Frobenius-equivariant maps from \(H^*(X)\) to \(H^*(Y_\nu)\). Extend
scalars to \(E\), sum, and project by \(e\). Every map to \(W=eR(T)\) is
obtained by projecting its composition with the inclusion
\(W\hookrightarrow R(T)\), so (CC-XC.5) implies the surjectivity of
(CC-XC.1). This remains a condition on a specified finite family and
specified codimension. It requires no Tate assertion for all varieties
or all codimensions. The projected condition can be weaker than full
surjectivity in (CC-XC.5).

Without semisimplicity, surjectivity onto
\(\operatorname{Hom}_{E[F]}(V,W)\) alone would not guarantee an
invertible map: all such maps could have nontrivial kernel. The
determinant criterion is the exact formulation retaining that issue.

\subsection{CC-XC04 --- Tate and numerical--homological equality after
the
criterion}\label{cc-xc04-tate-and-numericalhomological-equality-after-the-criterion}

\textbf{Corollary.} Suppose the equivalent conditions of CC-XC02 hold
for \(X\). Let \(k'/k\) be any finite extension, let \(n\ge0\), let
\(A/k'\) be an abelian variety and \(S/k'\) finite étale, and put

\[
 Z=X_{k'}^{\,n}\times A\times S.
\]

For the fixed prime \(\ell\), Frobenius on
\(H^*(Z_{\bar k},\mathbf Q_\ell)\) is semisimple; the rational Tate
cycle map is surjective in every codimension; and numerical equivalence
and \(\ell\)-adic homological equivalence agree on rational algebraic
cycles of \(Z\).

\textbf{Complete proof.} If \(Z\) is empty, all assertions are vacuous,
so assume it is nonempty. Base change of algebraic correspondences
preserves the isomorphism \(h_{\ell,E}(X)\simeq M\) proved in CC-XC02.
Its base change remains an abelian--Artin motive over \(k'\). Tensoring,
and then multiplying by \(h(A)\) and \(h(S)\), gives an isomorphism of
\(h_{\ell,E}(Z)\) with an object of \(\mathcal A_{q',E}\), where
\(q'=\#k'\). AQ.15 applies over every finite field. It proves
semisimplicity after extension to \(E\), and semisimplicity over
\(\mathbf Q_\ell\) follows because a characteristic-zero operator is
semisimple if and only if its extension to an algebraic closure is.

For codimension \(r\), fullness on the abelian sector identifies
algebraic morphisms \(\mathbf 1\to h_{\ell,E}(Z)(r)\) with the
Frobenius-invariant vectors of its realization. The weight decomposition
shows that only cohomological degree \(2r\) contributes. Consequently

\[
 \operatorname{im}\!\left(
 \mathrm{CH}^r(Z)_{\mathbf Q}\otimes_{\mathbf Q}E
 \longrightarrow H^{2r}(Z_{\bar k},E(r))
 \right)
 =H^{2r}(Z_{\bar k},E(r))^{F=1}.                 \tag{CC-XC.6}
\]

Here is the descent to the usual \(\mathbf Q_\ell\)-valued Tate
statement. Let \(C_\ell\) denote the image of
\(\mathrm{CH}^r(Z)_{\mathbf Q}\otimes_{\mathbf Q}\mathbf Q_\ell\) in
\(U=H^{2r}(Z_{\bar k},\mathbf Q_\ell(r))^{F=1}\). Flat scalar extension
commutes both with that image and with the kernel defining \(F=1\)
invariants. Equation (CC-XC.6) therefore says
\(C_\ell\otimes_{\mathbf Q_\ell}E=U\otimes_{\mathbf Q_\ell}E\). Faithful
flatness of a field extension gives \(U/C_\ell=0\). This argument uses
the image subspace; it does not presume injectivity after scalar
extension of rational homological classes.

Write \(D=\dim Z\). The Poincaré pairing

\[
 H^{2r}(Z_{\bar k},\mathbf Q_\ell(r))
 \times
 H^{2D-2r}(Z_{\bar k},\mathbf Q_\ell(D-r))
 \longrightarrow\mathbf Q_\ell
\]

is perfect and Frobenius invariant. Its restriction to the two invariant
subspaces is perfect: after diagonalizing semisimple Frobenius over
\(E\), eigenvalues \(\lambda,\mu\) can pair nontrivially only when
\(\lambda\mu=1\); the \(1\)-eigenspace is therefore paired perfectly
with the other \(1\)-eigenspace. Nondegeneracy descends to
\(\mathbf Q_\ell\).

If a rational cycle \(z\) is numerically trivial but
\(a=\operatorname{cl}_\ell(z)\ne0\), this invariant perfect pairing
provides an invariant complementary class \(b\) with
\(\langle a,b\rangle\ne0\). The already proved rational Tate
surjectivity expresses \(b\) as a finite \(\mathbf Q_\ell\)-linear
combination of classes of rational complementary cycles. Numerical
triviality makes every corresponding intersection degree zero,
contradicting the chosen pairing. Thus numerical triviality implies
homological triviality. The converse follows directly from compatibility
of the cycle class with intersection degrees. This proves the assertion
on rational cycles without assuming numerical realization at the start.
\(\square\)

The conclusion is for the specified prime \(\ell\) and this class of
products and base extensions. It does not assert the determinant
condition for arbitrary \(X\), or a comparison with another prime
without an additional argument.

\subsection{Relation to Milne and the remaining
question}\label{relation-to-milne-and-the-remaining-question}

Milne (1994), \textbf{Proposition 1.15 and Corollary 1.16},
author-manuscript pp.~11--12, derive semisimplicity,
numerical--homological equality and full realization from the full Tate
conjecture. \textbf{Proposition 2.6 and Remark 2.7}, p.~15, then give
the global abelian/Artin-generation conclusion under that hypothesis.
Those inspected statements do not themselves state the particular
one-sided finite-target criterion above. Here we prove that criterion
using only the already settled abelian-sector realization and the
specified cycle condition for \(X\). No historical priority claim is
made.

For a particular \(X\), the unresolved requirement has therefore been
made concrete: construct finitely many cross-correspondences whose
projected realization has nonzero determinant. Pure numerical Frobenius
eigenvalues alone do not establish this. The adopted RH and Dirichlet
GRH do not provide these cycles or prove the determinant condition.

\subsection{Sources and locators}\label{sources-and-locators}

\begin{itemize}
\tightlist
\item
  Deligne, P. (1974). La conjecture de Weil: I. \emph{Publications
  Mathématiques de l'IHÉS, 43}, 273--307. \textbf{Theorem 1.6, p.~276}.
  https://www.numdam.org/item/PMIHES\_1974\_\_43\_\_273\_0.pdf.
\item
  Honda, T. (1968). Isogeny classes of abelian varieties over finite
  fields. \emph{Journal of the Mathematical Society of Japan, 20}(1--2),
  83--95. https://doi.org/10.2969/jmsj/02010083. Main Theorem, p.~83;
  primary PDF:
  https://people.math.rochester.edu/faculty/doug/otherpapers/Honda-Ab.pdf.
\item
  Milne, J. S. (1994). Motives over finite fields. In U. Jannsen, S. L.
  Kleiman, \& J.-P. Serre (Eds.), \emph{Motives} (Proceedings of
  Symposia in Pure Mathematics, Vol. 55, Part 1, pp.~401--459). American
  Mathematical Society. https://jmilne.org/math/articles/1994aP.pdf.
  Author-manuscript pp.~5--6 for finite-field projectors, corrected
  symmetry and duality; pp.~11--12 for Proposition 1.15 and Corollary
  1.16; p.~15 for Proposition 2.6 and Remark 2.7.
\item
  Proved internal inputs: \textbf{CC-AG01} in
  \nolinkurl{research/consequences_criteria_abelian_group.tex};
  \textbf{AQ.15 / AQ-ABELIAN-RECOGNITION} and \textbf{AQ-WEIL-EXIST} in
  \nolinkurl{research/archive_questions.tex}. Their exact all-abelian primary
  inputs and limits are recorded in CC-TATE01--02 and the corresponding
  source ledger.
\end{itemize}
